% 第七章 系统测试与效果分析

\chapter{系统测试与效果分析}


\section{验证范围与实验环境}

FreeFlow 这一系统，它同时涉及到了桌面端、后端、智能合约，还有外部的 IPFS 服务。鉴于一部分功能需要依赖桌面端的本地代理、外部存储服务以及用户的交互流程，本文把验证工作划分成了两类能够稳定复现的任务：第一类，是凭借 Hardhat 本地链来开展的智能合约自动化测试，它的作用是去验证发布、购买还有访问控制等链上核心逻辑；第二类，则是基于离线评估集所进行的检索排序实验，用来对第 6 章所提出的资源发现策略进行验证。像登录、素材上传、详情展示、播放和评论等这些桌面端与后端相协同的功能，在第五章当中已经以业务链路的形式，对它们的实现和联调结果做出了说明，所以本章的实验验证，会把焦点聚集在前面提到的那两类可以重复执行的部分上。

为了保障实验过程能够复现，本文将检索排序评估和智能合约测试这两者，整合为了一套统一的实验流程，并且把结果以结构化的数据形式保存下来，专门用于本章统计图表的生成与分析工作。

实验所使用的硬件、软件与数据规模配置如表\ref{tab:test_environment}所示。

\begin{table}[htbp]
    \linespread{1.5}
    \zihao{5}
    \songti
    \centering
    \caption{实验环境配置}\label{tab:test_environment}
    \begin{tabular}{|l|p{10cm}|}
        \hline
        类别      & 环境说明                           \\ \hline
        硬件设备    & Apple M4，16 GB 内存              \\ \hline
        操作系统    & macOS 26.1（arm64）              \\ \hline
        后端运行时   & Node.js v20.18.0，pnpm 10.0.0   \\ \hline
        桌面端技术栈  & Electron、React、TypeScript      \\ \hline
        后端技术栈   & Node.js、Express、Prisma         \\ \hline
        合约测试工具链 & Hardhat 2.28.6，Solidity 0.8.26 \\ \hline
        排序评估集规模 & 56条资源，62个查询                    \\ \hline
        参考时间点   & 2026-04-22 12:00:00 UTC        \\ \hline
    \end{tabular}
\end{table}


\section{智能合约自动化测试结果}

智能合约测试基于Hardhat测试框架完成，围绕作品发布、购买授权、访问控制和销售配置修改等核心场景设计自动化用例。测试总计4项，全部通过，单次执行耗时不到1 s。测试结果如表\ref{tab:contract_tests}所示。

\begin{table}[htbp]
    \linespread{1.45}
    \zihao{5}
    \songti
    \centering
    \caption{智能合约自动化测试结果}\label{tab:contract_tests}
    \begin{tabular}{|l|p{8.3cm}|l|}
        \hline
        编号  & 测试内容                            & 结果 \\ \hline
        C01 & 付费作品发布后，用户支付正确金额可获得ERC-1155访问凭证 & 通过 \\ \hline
        C02 & 公开作品发布后，未购买用户默认拥有访问权限           & 通过 \\ \hline
        C03 & ERC-1155访问凭证不可在普通用户之间转让         & 通过 \\ \hline
        C04 & 只有创作者本人可以修改作品销售配置               & 通过 \\ \hline
    \end{tabular}
\end{table}

该结果表明，\texttt{PlatformHub}与\texttt{MusicAccess1155}之间的协作逻辑能够正确支撑作品发布、购买和访问校验；同时，创作者侧收益分配能够按预设流程导入\texttt{RoyaltySplitter}。这为桌面端调用链上接口提供了可靠前提。生产环境安全审计和大规模链上压力测试仍需在后续工作中继续完成。


\section{检索排序实验设置}

离线检索评估集是围绕第6章所提出的链上资源排序目标来进行构建的。这个评估集包含了56条结构化的模拟资源，并且按照14个艺术家/专辑簇来开展组织工作，它覆盖了标题、艺术家、专辑、流派、Token ID、合约地址、CID、resourceKey以及交易哈希等等诸多字段；同时，62个查询会去覆盖精确检索、模糊检索、链上身份定位以及输入扰动等场景。针对标题、Token ID、CID以及resourceKey这类单目标查询，仅仅会标注一个三级相关结果；而对于艺术家、专辑和流派查询，就会按照人工所标注的3/2/1相关等级来对排序质量进行评价。

这里需要指出的是，这个评估集主要是为当前系统在冷启动场景下开展可复现的实验来提供服务的。它虽然覆盖了当前实现所关心的查询类型，但并没有去模拟真实线上平台的全部流量特征。为了便于调试以及进行解释，部分链上身份字段选用了结构化测试标识，所以实验结果更适合用来说明排序机制在当前实现当中的工程有效性。而生产环境里的最终搜索质量，仍然需要凭借真实的日志来开展进一步的验证工作。

\begin{table}[htbp]
    \linespread{1.45}
    \zihao{5}
    \songti
    \centering
    \caption{检索评估查询类型分布}\label{tab:ranking_query_categories}
    \begin{tabular}{|l|c|l|}
        \hline
        查询类型           & 数量 & 说明                                            \\ \hline
        标题精确匹配         & 12 & 直接定位单首作品                                      \\ \hline
        艺术家查询          & 10 & 返回同创作者多首作品                                    \\ \hline
        专辑查询           & 10 & 验证同专辑资源排序                                     \\ \hline
        流派查询           & 8  & 验证分组标签下的多结果排序                                 \\ \hline
        Token ID 查询    & 3  & 验证链上标识精确定位                                    \\ \hline
        合约地址查询         & 2  & 验证MusicAccess合约检索                             \\ \hline
        CID 查询         & 2  & 验证IPFS内容寻址检索                                  \\ \hline
        resourceKey 查询 & 2  & 验证链上资源键精确定位                                   \\ \hline
        交易哈希查询         & 1  & 验证发布交易定位                                      \\ \hline
        拼写扰动查询         & 5  & 如\texttt{The Fields of Ard Skellix}           \\ \hline
        紧凑输入/空格扰动      & 4  & 如\texttt{TheFieldsofArdSkellig}、\texttt{新 地球} \\ \hline
        全角输入           & 1  & 验证NFKC归一化                                     \\ \hline
        描述关键词查询        & 1  & 验证长文本辅助召回                                     \\ \hline
        多词组合查询         & 1  & 验证艺术家与专辑联合意图                                  \\ \hline
    \end{tabular}
\end{table}

对比方法一共设置了4组。其中，数据库字段包含匹配基线这个方法只依赖于字段包含关系；仅文本评分方法则仅仅去计算标题、艺术家、专辑、流派、描述以及模糊相似度；文本+链上身份字段方法会在前面方法的基础上，增加Token ID、合约地址、CID、resourceKey以及交易哈希等等身份特征；而完整的排序算法也就是本文所实现的\texttt{chain\_resource\_rank\_v1}，它还会进一步引入新鲜度以及热度信号。评价指标方面，囊括了Precision@10、MRR、nDCG@10、零结果率以及平均响应时间。


\section{检索效果对比实验}

表\ref{tab:ranking_metrics_actual}和图\ref{fig:ranking_metrics_bar}给出了4种方法在62个查询上的总体表现。由于评估集包含较多单目标定位查询，Precision@10的绝对值受到上限约束，因此本节重点结合MRR和nDCG@10分析排序质量。

\begin{table}[htbp]
    \linespread{1.45}
    \zihao{5}
    \songti
    \centering
    \caption{检索排序指标对比结果}\label{tab:ranking_metrics_actual}
    \begin{tabular}{|p{3.2cm}|c|c|c|c|c|}
        \hline
        方法          & Precision@10 & MRR    & nDCG@10 & 零结果率    & 平均响应时间/ms \\ \hline
        数据库字段包含匹配基线 & 0.1790       & 0.8468 & 0.8489  & 12.90\% & 0.2350    \\ \hline
        仅文本评分       & 0.1806       & 0.8306 & 0.8237  & 16.13\% & 3.1300    \\ \hline
        文本+链上身份字段   & 0.1952       & 0.9758 & 0.9646  & 0.00\%  & 3.1406    \\ \hline
        完整排序算法      & 0.1952       & 0.9758 & 0.9646  & 0.00\%  & 3.3014    \\ \hline
    \end{tabular}
\end{table}

\begin{figure}[htbp]
    \centering
    \begin{tikzpicture}
        \begin{axis}[
            ybar,
            width=0.98\textwidth,
            height=8.1cm,
            bar width=9pt,
            ymin=0,
            ymax=105,
            ylabel={指标值（\%）},
            symbolic x coords={sql_contains,text_only,text_identity,full_ranker},
            xtick=data,
            xticklabels={字段包含匹配基线,仅文本评分,文本+链上身份字段,完整排序算法},
            xticklabel style={rotate=18,anchor=east,font=\zihao{6}},
            ymajorgrids=true,
            legend style={at={(0.5,1.05)},anchor=south,legend columns=3,font=\zihao{6}},
            enlarge x limits=0.08,
        ]
            \addplot table[x=method_key,y=precision_at_10_pct,col sep=comma]{figures/data/search_ranking_metrics.csv};
            \addplot table[x=method_key,y=mrr_pct,col sep=comma]{figures/data/search_ranking_metrics.csv};
            \addplot table[x=method_key,y=ndcg_at_10_pct,col sep=comma]{figures/data/search_ranking_metrics.csv};
            \legend{Precision@10,MRR,nDCG@10}
        \end{axis}
    \end{tikzpicture}
    \caption{不同排序策略的检索质量对比}
    \label{fig:ranking_metrics_bar}
\end{figure}

从实验结果可以看出，单纯依赖数据库字段包含匹配虽然开销最低，但在拼写扰动、空格扰动和链上复杂身份查询场景中会出现明显漏召回，零结果率达到 12.90\%。仅文本评分方法能够处理一部分模糊查询，但由于缺少身份字段特征，对Token ID、合约地址、CID 和 resourceKey 等链上定位查询支持不足，其 MRR 和 nDCG@10 分别只有 0.8306 和 0.8237，且零结果率进一步上升到 16.13\%。

当加入链上身份字段后，零结果率下降到 0，MRR 从 0.8306 提升到 0.9758，nDCG@10 提升到 0.9646，说明绝大多数查询都能在首位或前两位命中目标资源。完整排序算法在此基础上继续引入新鲜度和热度特征，在 Precision@10、MRR 和 nDCG@10 上保持与文本+ 链上身份字段方法一致的总体表现，同时为在线检索结果补充了可解释因子，说明新增排序特征没有破坏结果稳定性。


\section{典型查询案例分析}

完整排序算法除了输出结果列表外，还会返回排序原因，便于分析命中特征与排序依据。表\ref{tab:ranking_case_analysis}给出了3类典型查询的实验结果。

\begin{table}[htbp]
    \linespread{1.45}
    \zihao{5}
    \songti
    \centering
    \caption{典型查询案例分析}\label{tab:ranking_case_analysis}
    \begin{tabular}{|p{3.6cm}|p{3.6cm}|p{5.2cm}|}
        \hline
        查询输入                               & 首位结果                      & 主要排序原因                          \\ \hline
        \texttt{The Fields of Ard Skellix} & The Fields of Ard Skellig & 标题分词命中80\%，文本相似度95\%，能够纠正轻微拼写错误 \\ \hline
        \texttt{新 地球}                      & 新地球                       & 标题经紧凑化后精确匹配，能够处理空格扰动输入          \\ \hline
        resourceKey精确查询                    & White Palace              & resourceKey精确匹配，同时命中音乐合约地址      \\ \hline
    \end{tabular}
\end{table}

上述结果表明，排序模块不仅能够处理传统文本检索，还能处理链上身份字段定位以及带空格、大小写和轻微拼写误差的输入。排序原因字段使系统具备较强的可解释性，也便于后续继续调权和排查误召回问题。


\section{性能测试结果}

为了观察完整排序算法在不同候选规模下的开销变化，本文将候选集合分别扩展到50、100、300和500条，并统计平均响应时间与95\%分位响应时间。结果如表\ref{tab:ranking_latency}和图\ref{fig:ranking_latency_curve}所示。

\begin{table}[htbp]
    \linespread{1.45}
    \zihao{5}
    \songti
    \centering
    \caption{不同候选规模下的排序耗时}\label{tab:ranking_latency}
    \begin{tabular}{|c|c|c|}
        \hline
        候选数量 & 平均响应时间/ms & 95\%分位响应时间/ms \\ \hline
        50   & 3.80      & 6.73          \\ \hline
        100  & 7.49      & 13.42         \\ \hline
        300  & 21.99     & 39.68         \\ \hline
        500  & 36.09     & 64.37         \\ \hline
    \end{tabular}
\end{table}

\begin{figure}[htbp]
    \centering
    \begin{tikzpicture}
        \begin{axis}[
            width=0.94\textwidth,
            height=7.8cm,
            xlabel={候选数量},
            ylabel={响应时间/ms},
            xmin=40,
            xmax=510,
            ymin=0,
            ymax=90,
            ymajorgrids=true,
            legend style={at={(0.5,1.05)},anchor=south,legend columns=2,font=\zihao{6}},
            tick label style={font=\zihao{6}},
            label style={font=\zihao{6}},
        ]
            \addplot+[mark=*,line width=1pt] table[x=candidate_count,y=avg_latency_ms,col sep=comma]{figures/data/search_latency_candidates.csv};
            \addplot+[mark=square*,line width=1pt] table[x=candidate_count,y=p95_latency_ms,col sep=comma]{figures/data/search_latency_candidates.csv};
            \legend{平均响应时间,95\%分位响应时间}
        \end{axis}
    \end{tikzpicture}
    \caption{完整排序算法在不同候选规模下的延迟变化}
    \label{fig:ranking_latency_curve}
\end{figure}

结果显示，随着候选数量增加，平均响应时间近似线性增长，这与第 6 章中的复杂度分析一致。在 500 条候选资源下，完整排序算法的平均响应时间为 36.09 ms，95\% 分位响应时间为 64.37 ms，仍处于桌面端检索交互可接受范围内。因此，在本文面向的系统规模与应用场景下，\texttt{chain\_resource\_rank\_v1}能够在保证排序质量的同时维持可接受的在线响应开销。


\section{结果讨论与有效性说明}

综合本章结果，可以得出两点较为稳妥的结论。第一，系统当前已经在可复现范围内验证了链上发布、购买和访问控制这条关键业务链路，说明合约层设计能够支撑作品凭证化和访问授权。第二，检索排序模块在冷启动场景下具备较好的工程可用性，加入链上身份字段后可以显著降低零结果率，并维持较低延迟。

同时，本章也有明确边界。桌面端的登录、上传、详情展示、播放、评论和下载等功能已经完成实现与联调，但尚未形成大规模自动化端到端测试。检索实验所使用的评估集是结构化模拟数据，适合验证算法链路和特征设计是否工作正常，但不足以替代真实用户搜索日志和线上A/B测试。因此，本文的结论应表述为：FreeFlow完成了面向音乐场景的混合架构实现验证，生产环境中的完整平台成熟度仍有待进一步提升。

